%======================================================================
\section{Validity theorems}
\label{sec:validity}

This section proves that validity is structural: every admissible
coefficient vector prices from a genuine mean-one law, so the optimizer of
\cref{sec:calibration} never sees a constraint.  ``Admissible'' means
exactly one thing:
\begin{equation}
\Theta_N=\big\{\theta=(L,R,a_2,\dots,a_N)\in\R^{N+1}:\ \lambda_+<1\big\},
\label{eq:admissible}
\end{equation}
an \emph{open} half-space-like subset of $\R^{N+1}$ cut out by one smooth
scalar inequality --- and \cref{sec:computation} shows a change of variable
under which even that inequality disappears, the admissible set becoming
all of $\R^{N+1}$.

\subsection{One expiry: butterfly-freedom without a certificate}

\begin{proposition}[One-expiry no-arbitrage]
\label{prop:noarb}
Let $\theta\in\Theta_N$ and let $c(k)=\E[(e^{X}-e^{k})^{+}]$ be priced from
the slice \eqref{eq:lqd}.  Then:
\begin{enumerate}[leftmargin=1.6em,label=(\roman*)]
\item $X$ has a strictly positive, continuous density on $\R$, and
$Y=e^{X}$ a strictly positive density on $(0,\infty)$ with $\E[Y]=1$;
\item as a function of the strike $y=e^{k}$, the call is strictly
decreasing and strictly convex, with $(1-y)^{+}\le c\le1$, slope
$c'(y)=-\Prob(Y>y)\in(-1,0)$, and $c\to0$ as $y\to\infty$;
\item every finite-width butterfly has strictly positive price, and
put--call parity $c(k)-P(k)=1-e^{k}$ holds identically;
\item convexity in \emph{log}-strike fails where it should: $c$ is concave
in $k$ on a deep in-the-money region, and this is a property of every
arbitrage-free slice, not a defect.
\end{enumerate}
\end{proposition}

\begin{proof}
(i) is \eqref{eq:density} together with \cref{prop:shift}: $x$ is a
$C^1$ diffeomorphism of $\R$ (its derivative $e^{g(\Lambda(z))}$ is
continuous and positive), so $X=x(Z)$ has density
$\rho(x^{-1}(\cdot))/x'(x^{-1}(\cdot))>0$; the shift $m$ makes
$\E[e^{X}]=1$, and $Y=e^{X}$ inherits a positive density on $(0,\infty)$
by another smooth change of variable.

(ii) Write the call rank-by-rank: with $U$ uniform and $Y=e^{Q(\logit U)}$
--- hereafter we write $Q(u)$ for $x(\logit u)$ ---
\[
c(y)=\int_0^1\big(e^{Q(u)}-y\big)^{+}\dd u .
\]
For fixed $u$ the integrand is convex and non-increasing in $y$; both
properties survive the integral.  Strictness and the slope value follow
from differentiating under the integral (dominated by $1$):
$c'(y)=-\int_0^1\1\{e^{Q(u)}>y\}\dd u=-\Prob(Y>y)$, which lies strictly
between $-1$ and $0$ because the density of $Y$ is positive everywhere,
and is strictly increasing in $y$ for the same reason, which is strict
convexity.  The bounds are $\E[(Y-y)^{+}]\le\E[Y]=1$ and
$(1-y)^{+}=(\E[Y]-y)^{+}\le\E[(Y-y)^{+}]$ by Jensen; $c(y)\le
\E[Y\1\{Y>y\}]\to0$ by dominated convergence.

(iii) A butterfly's price is a second difference of a convex function,
positive by (ii); strictness because convexity is strict.  Parity:
$c(k)-P(k)=\E[(Y-e^{k})^{+}-(e^{k}-Y)^{+}]=\E[Y-e^{k}]=1-e^{k}$.

(iv) In log-strike, $\partial^2 c/\partial k^2
=y^2\,\partial^2c/\partial y^2+y\,\partial c/\partial y
=e^{k}\big(e^{k}f_Y(e^{k})-\Prob(Y>e^{k})\big)$.  As $k\to-\infty$,
$e^{k}f_Y(e^{k})=f_X(k)\to0$ while $\Prob(Y>e^{k})\to1$, so the bracket is
eventually negative: $c$ is concave in $k$ deep in the money.  The
computation used only the existence of a density of $X$ vanishing at
$-\infty$, so this holds for essentially every arbitrage-free slice in the
sense of \cref{def:slice}, not specially for LQD.
\end{proof}

Point (iv) is recorded because it once cost a real audit an afternoon: a
convexity check written in log-strike will flag every correct slice as
arbitrageable in the deep left wing.  Validity checks belong in the strike
variable $y$, where \cref{def:slice} lives.  \Cref{fig:butterfly}
illustrates (ii) and (iii) on a fitted real-market slice: the call curve
sits below its secant chords everywhere, and finite-width butterflies
priced off the fitted slice recover the model density as their width
shrinks --- \eqref{eq:BL} made visible.

\begin{figure}[htbp]
\centering
\paperfig{\textwidth}{fig_butterfly.pdf}
\caption{Validity as geometry, on a fitted real-market slice (SPY, December
2026, from the frozen snapshot).  Panel A: the normalized call curve
against strike $y=e^k$ with a secant chord drawn --- the curve is convex,
every chord sits above it, so every butterfly has non-negative price.
Panel B: discrete butterflies of fixed half-width, priced off the fitted
slice and divided by the squared half-width, against the model's own
density (line): the Breeden--Litzenberger identity \eqref{eq:BL} realized
by actual portfolios, to a worst relative gap of
\MacButterflyDensityRelErrPct\%.  Butterfly-freedom is not something this
fit achieved; it is something the representation cannot lose.}
\label{fig:butterfly}
\end{figure}

\subsection{The wall is a fact, not a preference}

\begin{proposition}[Finite-forward wall]
\label{prop:wall}
Within the family, the normalizing integral converges --- and hence the
slice exists --- if and only if $\lambda_+<1$ (\cref{prop:shift}).
Moreover the wall is model-free: \emph{any} law of $X$ whose right tail
satisfies $\Prob(X>x)\ge\kappa e^{-x}$ for some $\kappa>0$ and all large
$x$ has $\E[e^{X}]=\infty$.  No parametrization can price a forward from
such a law.
\end{proposition}

\begin{proof}
The first claim is \cref{prop:shift}.  For the second, integrate by parts,
or directly: for any level $x_0$,
$\E[e^{X}]\ge\E[e^{X}\1\{X>x_0\}]
=\int_0^\infty \Prob(e^{X}\1\{X>x_0\}>t)\,\dd t
\ge\int_{e^{x_0}}^{\infty}\Prob(X>\log t)\,\dd t
\ge\kappa\int_{e^{x_0}}^{\infty}t^{-1}\,\dd t=\infty$.
\end{proof}

In words: $\lambda_+$ is the exponential scale of the right return tail
(\cref{sec:tails}), and a right tail at scale $\ge1$ in $X$ is a return
distribution whose mean does not exist.  The wall $\lambda_+<1$ is the
price of a finite forward, charged by probability theory, not by the
model.  There is no corresponding left-hand wall: $\lambda_-$ may be
arbitrarily large, because $e^{X}\le1$ on the left half of the law's
support contributes at most $1$ to the mean.

\subsection{What the coordinates cover, and what they do not}

The title of this section would be overclaimed by ``coordinates for the
space of smiles''.  The honest statement is a density result \emph{within
a class}, with the exclusions named.

\begin{proposition}[Universality, with exclusions]
\label{prop:universal}
Let $X^{\ast}$ be a mean-one law whose quantile function $Q^{\ast}$ is
differentiable with
\[
g^{\ast}(u)\;=\;\log Q^{\ast\prime}(u)+\log u+\log(1-u)
\]
extending continuously to $[0,1]$ --- equivalently, a law with a positive
continuous density and exponential tails, of scales
$\lambda_-^{\ast}=e^{g^{\ast}(0)}>0$ and $\lambda_+^{\ast}=e^{g^{\ast}(1)}<1$.
Then:
\begin{enumerate}[leftmargin=1.6em,label=(\roman*)]
\item (\emph{approximation}) there are admissible LQD slices of increasing
order $N$ whose $g_N\to g^{\ast}$ uniformly on $[0,1]$; their quantiles
converge to $Q^{\ast}$ uniformly on compact subsets of $(0,1)$, their
normalized call curves converge uniformly in $k$, and the laws of the
gross returns converge in Wasserstein-$p$ distance for every
$p\in[1,\,1/\lambda_+^{\ast})$ --- for laws on $\R$ this metric has the
concrete form
$W_p(\mu,\nu)^{p}=\int_0^1|Q_\mu(u)-Q_\nu(u)|^{p}\,\dd u$, the $p$-th
mean cost of carrying one law onto the other rank by rank, which is why
it is the natural metric for a quantile-side model;
\item (\emph{exclusions}) a law with an atom, an interval of zero density
inside its support, bounded support, point mass at zero (default), or a
Gaussian or faster-decaying tail is not in the family at any finite
order.  Moreover such laws are unreachable in the limit, in the
following precise sense.  Let admissible slices have log-speeds $g_N$
converging uniformly on $[0,1]$; the limit $g_\infty$ is then continuous
with $g_\infty(1)\le0$, and two branches exhaust the possibilities.  If
$g_\infty(1)<0$, the laws converge weakly and the limit belongs to the
class of (i).  If $g_\infty(1)=0$, weak convergence can fail outright
--- the martingale shifts can diverge and the mass escape to $-\infty$
--- and \emph{if} the laws do converge weakly, the limit is a boundary
law at the wall: it still has a strictly positive continuous density on
all of $\R$ (no atom, no gap, no bounded support, no
Gaussian-or-faster tail); what is lost is only the tail scale --- its
moment strip collapses to $r\le1$, and by Fatou it need not even retain
unit mean.  In neither branch is an excluded law reached.
\end{enumerate}
\end{proposition}

The proof is in \cref{app:proofs}; the mechanism deserves three sentences
here.  Approximation works because admissibility is an \emph{open}
condition: a Weierstrass polynomial approximation $g_N$ of $g^{\ast}$ has
$g_N(1)\to g^{\ast}(1)<0$, so $g_N$ is itself admissible for $N$ large ---
the family is stable under the approximation that fills it out, which is
precisely what a closed condition (an equality constraint, say) would
ruin.  The exclusions are equally structural: an atom is a jump of the
quantile function, a gap is a flat of it, and both are impossible for
$Q=x\circ\logit$ with $x'$ a positive continuous exponential of a
polynomial, at any order; a Gaussian tail forces $g^{\ast}(u)\to-\infty$
at the endpoint, off the chart entirely.

For desk use the exclusions matter as approximation targets --- the
double-hump example of \cref{sec:examples} fits a genuinely bimodal
density and measures what a finite $N$ delivers --- and as honest limits:
a pinned settlement print (an atom) or a hard default mass will be
\emph{smoothed}, not represented, and the fitter should say so rather than
pretend otherwise.  \Cref{sec:limits} returns to this contract.

The scope sentence of the paper, then, earned rather than asserted: the
LQD construction provides \emph{unconstrained coordinates for a flexible
exponential-tail class of mean-one laws} --- all of that class, nothing
outside it, and the one inequality that delimits it is removed by the
logistic chart of \cref{sec:computation}.
